% =====================================================================
\section{Rebotes contra superficies fijas}
\label{sec:rebotes}
% =====================================================================

\begin{proposicion}[Rebote contra un plano fijo liso]
\label{prop:rebote}
Para una partícula que choca contra un plano fijo liso (límite $m_{2}\to\infty$, $v_{2}=0$):
\begin{equation}
  \resultado{v_{n}'=-e\,v_{n},\qquad v_{t}'=v_{t}},\qquad
  Q=\tfrac12\,m\,(1-e^{2})\,v_{n}^{2} .
  \label{eq:rebote}
\end{equation}
\end{proposicion}

\begin{proof}
En \eqref{eq:frontal-explicito} con $v_{2}=0$ y $m_{1}/m_{2}\to0$: $V\cm\to0$, luego $v_{1n}'=-e\,v_{1n}$; $\mu\to m_{1}$
en $Q$. El plano liso sólo ejerce percusión normal: $v_{t}$ no cambia (\cref{prop:lisas}).
\end{proof}

\begin{corolario}[Pelota que bota]
\label{cor:bote}
Soltada desde $h$ (sin rozamiento con el aire), tras el $n$-ésimo bote
\begin{equation}
  v_{n}=e^{n}\sqrt{2gh},\qquad h_{n}=e^{2n}h,\qquad
  \resultado{e=\sqrt{h_{n+1}/h_{n}}} ,
\end{equation}
y el tiempo total hasta el reposo es finito:
$T_{\infty}=\sqrt{\tfrac{2h}{g}}\Bigl(1+2\sum_{n\ge1}e^{n}\Bigr)=\sqrt{\tfrac{2h}{g}}\,\dfrac{1+e}{1-e}$.
\end{corolario}

\begin{proof}
Por \eqref{eq:rebote}, $v_{n}=e\,v_{n-1}$; por energía, $h_{n}=v_{n}^{2}/2g$. Cada vuelo dura
$2v_{n}/g=2e^{n}\sqrt{2h/g}$; la serie geométrica converge porque $e<1$.
\end{proof}

% ---------------------------------------------------------------------
\begin{figure}[ht]
\centering
\begin{tikzpicture}
\begin{axis}[oscuro,width=12cm,height=5.0cm,xmin=0,xmax=9.4,ymin=0,ymax=1.12,
  xlabel={$t/\sqrt{2h/g}$},ylabel={$y/h$},ytick={0,0.41,0.64,1},
  yticklabels={$0$,$e^{4}$,$e^{2}$,$1$}]
  \addplot[cian,line width=0.9pt] table[x=t,y=y]{datos/bote.dat};
  \draw[aux] (axis cs:9,0) -- (axis cs:9,0.62);
  \node[peq,anchor=south,align=center] at (axis cs:9,0.62) {$T_{\infty}$};
  \draw[aux] (axis cs:0,0.64) -- (axis cs:2.6,0.64);
  \draw[aux] (axis cs:0,0.4096) -- (axis cs:5.6,0.4096);
\end{axis}
\end{tikzpicture}
\caption{Pelota que bota con $e=0{,}8$: alturas $e^{2n}h$ y duración de los vuelos
proporcional a $e^{n}$. Los infinitos botes ocurren en un tiempo finito,
$T_{\infty}=\frac{1+e}{1-e}\sqrt{2h/g}=9\sqrt{2h/g}$ (discontinua vertical).}
\label{fig:bote}
\end{figure}
% ---------------------------------------------------------------------

\begin{clave}{Método · rebotes}
\begin{enumerate}[leftmargin=1.4em]
  \item Separar el movimiento: \emph{vuelo} (cinemática o energía con las fuerzas que haya)
        y \emph{bote} instantáneo (\cref{prop:rebote}).
  \item En el bote, invertir y multiplicar por $e$ sólo la componente normal; la horizontal se conserva.
  \item Encadenar: la velocidad de salida de un bote es la inicial del vuelo siguiente.
\end{enumerate}
\end{clave}

% =====================================================================
\subsection{Problemas: rebotes}
% =====================================================================

\begin{problema}{1}[Caída con resistencia cuadrática]
\label{prob:1}
Una partícula de masa $m$ en reposo cae desde una altura $h$ con resistencia del aire
$F_{\mathrm{aire}}=kv^{2}$, $k>0$. Hallar (a) la velocidad en función del espacio recorrido y
(b) la altura máxima tras el primer rebote si $e=1$.
\end{problema}

\begin{resolucion}
Truco: con $w\defeq v^{2}$, $\;m\,\od{v}{t}=m\,v\,\od{v}{s}=\dfrac m2\,\od{w}{s}$, y la ecuación es lineal en $w$.
\begin{enumerate}
  \item[(a)] \emph{Bajada} ($s$ hacia abajo, $w(0)=0$): $\dfrac m2\od{w}{s}=mg-kw$. Separando,
  \begin{gather*}
    \int_{0}^{w}\frac{\dd w'}{mg-kw'}=\frac2m\int_{0}^{s}\dd s'
    \ \Longrightarrow\ -\frac1k\ln\frac{mg-kw}{mg}=\frac{2s}{m}\ \Longrightarrow\ w=\frac{mg}{k}\bigl(1-\e^{-2ks/m}\bigr),\\
    \resultado{v(s)=v_{L}\sqrt{1-\e^{-2ks/m}}},\qquad v_{L}\defeq\sqrt{\frac{mg}k} .
  \end{gather*}
  Límites: $s\to\infty$, $v\to v_{L}$ (velocidad límite); $k\to0$, $v^{2}\to2gs$.
  \item[(b)] \emph{Bote} con $e=1$: sale hacia arriba con $w_{h}=v^{2}(h)=v_{L}^{2}\bigl(1-\e^{-2kh/m}\bigr)$.
  \emph{Subida} ($y$ hacia arriba; peso y rozamiento hacia abajo): $\dfrac m2\od{w}{y}=-(mg+kw)$, luego
  \[
    \frac1k\ln\frac{mg+kw_{h}}{mg+kw}=\frac{2y}{m}
    \ \xrightarrow{\ w=0\ }\ 
    H=\frac{m}{2k}\ln\Bigl(1+\frac{w_{h}}{v_{L}^{2}}\Bigr)
    =\resultado{\frac{m}{2k}\ln\bigl(2-\e^{-2kh/m}\bigr)} .
  \]
  \item[(c)] \emph{Comprobación}: desarrollando en $k$, $H\simeq h-\dfrac{2k}{m}h^{2}<h$; aunque $e=1$, el aire
  disipa energía en la bajada y en la subida (\cref{fig:aire}).\qedhere
\end{enumerate}
\end{resolucion}

% ---------------------------------------------------------------------
\begin{figure}[ht]
\centering
\begin{tikzpicture}
\begin{groupplot}[group style={group size=2 by 1,horizontal sep=16mm},
  oscuro,width=7.2cm,height=5.6cm,xmin=0,xmax=3]
  \nextgroupplot[ymin=0,ymax=1.15,xlabel={$ks/m$},ylabel={$v/v_{L}$},
    title={(a) Velocidad de caída}]
    \addplot[cian,domain=0:3,samples=150] {sqrt(1-exp(-2*x))};
    \draw[aux] (axis cs:0,1) -- (axis cs:3,1);
    \node[peq,anchor=south east] at (axis cs:3,1) {velocidad límite};
  \nextgroupplot[ymin=0,ymax=1.08,xlabel={$kh/m$},ylabel={$H/h$},
    title={(b) Altura tras el bote ($e=1$)}]
    \addplot[rosa,domain=0.005:3,samples=150] {ln(2-exp(-2*x))/(2*x)};
    \addplot[tenue,dashed,line width=0.7pt,domain=0.2:3,samples=60] {ln(2)/(2*x)};
    \node[peq,anchor=south west] at (axis cs:1.6,0.24) {$\frac{\ln2}{2kh/m}$};
\end{groupplot}
\end{tikzpicture}
\caption{\Cref{prob:1}. (a) La velocidad tiende a $v_{L}=\sqrt{mg/k}$. (b) Fracción de
altura recuperada: $1$ sin aire y $\ln2\,/(2kh/m)$ si la caída es muy larga (discontinua).}
\label{fig:aire}
\end{figure}
% ---------------------------------------------------------------------

\begin{problema}{7}[Pelota por una escalera]
\label{prob:7}
Una pelota baja una escalera rebotando en el centro de cada escalón, de altura y
profundidad $h$. Tras cada bote se eleva hasta el nivel del escalón anterior. Obtener el
coeficiente de restitución.
\end{problema}

\begin{resolucion}
Escalones lisos: $u=v_{x}$ constante. El movimiento es \emph{periódico}: el mismo
vuelo se repite de escalón en escalón (\cref{fig:escalera}).
\begin{enumerate}
  \item \emph{Subida tras el bote}: llega a altura $h$ sobre el escalón en el que bota, luego
        la componente vertical de salida es $v_{\uparrow}=\sqrt{2gh}$.
  \item \emph{Bajada}: desde ese vértice (nivel del escalón anterior) hasta el siguiente escalón,
        que está $h$ por debajo del actual, cae $2h$: $v_{\downarrow}=\sqrt{2g\cdot2h}=2\sqrt{gh}$.
  \item \emph{Bote}: $v_{\uparrow}=e\,v_{\downarrow}$, luego
        \[
          e=\frac{\sqrt{2gh}}{2\sqrt{gh}}=\resultado{\frac1{\sqrt2}\simeq0{,}71} .
        \]
  \item \emph{Velocidad horizontal} (comprobación de que el movimiento es posible): el vuelo dura
        $t_{\uparrow}+t_{\downarrow}=\sqrt{\tfrac{2h}g}+2\sqrt{\tfrac hg}$ y avanza $h$:
        $u=\dfrac{\sqrt{gh}}{2+\sqrt2}$. La pelota pasa el borde del escalón a $0{,}96\,h$ sobre él.\qedhere
\end{enumerate}
\end{resolucion}

% ---------------------------------------------------------------------
\begin{figure}[ht]
\centering
\begin{tikzpicture}[scale=1.15,
  declare function={uu=1/(2+sqrt(2));}]
  % escalera
  \fill[tenue!22!fondo] (-1,1) -- (0,1) -- (0,0) -- (1,0) -- (1,-1) -- (2,-1) -- (2,-2) -- (3,-2) -- (3,-3) -- (4.6,-3) -- (4.6,-3.35) -- (-1,-3.35) -- cycle;
  \draw[texto!70,line width=0.6pt] (-1,1) -- (0,1) -- (0,0) -- (1,0) -- (1,-1) -- (2,-1) -- (2,-2) -- (3,-2) -- (3,-3) -- (4.6,-3);
  % trayectorias
  \foreach \n in {-1,0,1,2}{
    \draw[cian,line width=0.9pt,smooth,samples=50]
      plot[domain=0:{sqrt(2)+2}] ({\n+0.5+uu*\x},{-\n+sqrt(2)*\x-0.5*\x*\x});
    \fill[cian] ({\n+0.5},{-\n}) circle (0.06);
  }
  \fill[cian] (3.5,-3) circle (0.06);
  % nivel del escalón anterior
  \draw[ambar,dash pattern=on 2pt off 2pt,line width=0.5pt] (1,0) -- (2.35,0);
  \node[peq,text=ambar,anchor=west] at (2.4,0) {vértice: nivel del escalón anterior};
  % cotas
  \draw[texto!80,line width=0.45pt,|-|] (0.78,0) -- (0.78,-1) node[midway,anchor=east,etiq,inner sep=3pt] {$h$};
  \draw[texto!80,line width=0.45pt,|-|] (2,-2.22) -- (3,-2.22) node[midway,anchor=north,etiq,inner sep=2pt] {$h$};
  \node[peq,anchor=north,align=center] at (1.5,-1.1) {bote};
\end{tikzpicture}
\caption{\Cref{prob:7} (dibujado a escala con $e=1/\sqrt2$). Cada vuelo sube $h$ sobre el
escalón del bote, baja $2h$ hasta el siguiente y avanza $h$ en horizontal.}
\label{fig:escalera}
\end{figure}
% ---------------------------------------------------------------------

\begin{problema}{10}[Mate tras un bote]
\label{prob:10}
Un pívot lanza un balón oblicuamente hacia el suelo desde $h_{0}=2{,}0$\,m. Tras botar, otro
jugador situado a una distancia horizontal $D=3{,}0$\,m lo recoge en su altura máxima
$H=3{,}0$\,m. Si $e=0{,}80$, determinar la velocidad inicial. Suelo horizontal y liso; sin
rozamiento con el aire.
\end{problema}

\begin{resolucion}
Ejes: $x$ horizontal desde el pívot, $y$ vertical hacia arriba; $\vb{v}_{0}=(v_{x},-v_{0y})$ con $v_{0y}>0$; $g=9{,}8$\,m/s$^{2}$.
Se resuelve \emph{hacia atrás}, desde el dato final (\cref{fig:mate}).
\begin{enumerate}
  \item \emph{Tras el bote}: sube hasta $H$, luego sale con $v_{\uparrow}=\sqrt{2gH}=7{,}67$\,m/s y tarda
        $t_{2}=v_{\uparrow}/g=0{,}783$\,s.
  \item \emph{Bote}: $v_{\downarrow}=v_{\uparrow}/e=\sqrt{2gH}/e=9{,}59$\,m/s; $v_{x}$ no cambia (suelo liso).
  \item \emph{Primer vuelo}: $v_{\downarrow}^{2}=v_{0y}^{2}+2gh_{0}$, luego
        \[
          v_{0y}=\sqrt{\frac{2gH}{e^{2}}-2gh_{0}}=7{,}26\text{ m/s},\qquad
          t_{1}=\frac{v_{\downarrow}-v_{0y}}{g}=0{,}237\text{ s}.
        \]
  \item \emph{Horizontal}: $D=v_{x}(t_{1}+t_{2})$, luego $v_{x}=\dfrac{3{,}0}{1{,}020}=2{,}94$\,m/s.
  \item \emph{Resultado}:
        \[
          \resultado{v_{0}=\sqrt{v_{x}^{2}+v_{0y}^{2}}\simeq7{,}8\text{ m/s},\qquad
          \arctan\frac{v_{0y}}{v_{x}}\simeq68^{\circ}\text{ bajo la horizontal}} .
        \]
        El bote se produce a $x_{b}=v_{x}t_{1}=0{,}70$\,m del pívot.\qedhere
\end{enumerate}
\end{resolucion}

% ---------------------------------------------------------------------
\begin{figure}[ht]
\centering
\begin{tikzpicture}
\begin{axis}[oscuro,width=10.5cm,height=7.2cm,clip=false,
  xmin=-0.5,xmax=3.6,ymin=0,ymax=3.6,xlabel={$x$ (m)},ylabel={$y$ (m)},
  xtick={0,0.7,1,2,3},xticklabels={$0$,$x_{b}$,$1$,$2$,$3$},ytick={0,1,2,3}]
  \addplot[cian,line width=1.0pt,domain=0:0.23749,samples=40,variable=\t]
    ({2.9413*t},{2-7.2578*t-4.9*t^2});
  \addplot[cian,line width=1.0pt,domain=0:0.78246,samples=80,variable=\t]
    ({0.69853+2.9413*t},{7.6681*t-4.9*t^2});
  \draw[vec,ambar,line width=0.8pt] (axis cs:0,2) -- (axis cs:0.1765,1.5645)
    node[anchor=west,etiq,ambar,inner sep=3pt] {$\vb{v}_{0}$};
  \node[punto,fill=texto] at (axis cs:0,2) {};
  \node[punto,fill=rosa] at (axis cs:3,3) {};
  \node[punto,fill=cian] at (axis cs:0.6985,0) {};
  \draw[aux] (axis cs:3,0) -- (axis cs:3,3);
  \draw[aux] (axis cs:0,0) -- (axis cs:0,2);
  \node[peq,anchor=west,inner sep=4pt] at (axis cs:3,3) {mate};
  \node[etiq,anchor=east,inner sep=3pt] at (axis cs:0,1) {$h_{0}$};
  \node[etiq,anchor=west,inner sep=3pt] at (axis cs:3,1.5) {$H$};
  \draw[tenue,line width=0.45pt,{Stealth[length=3.5pt]}-{Stealth[length=3.5pt]}] (axis cs:0,3.35) -- (axis cs:3,3.35)
    node[midway,anchor=south,etiq,inner sep=2pt] {$D$};
\end{axis}
\end{tikzpicture}
\caption{\Cref{prob:10}: trayectoria calculada (a escala). El balón sale a
$68^{\circ}$ bajo la horizontal, bota a $0{,}70$\,m y alcanza su altura máxima $H=3$\,m
justo en $x=D$.}
\label{fig:mate}
\end{figure}
% ---------------------------------------------------------------------
